% ==============================================================================
% MỤC 1.1 - 1.2: ĐẶT VẤN ĐỀ VÀ MỤC TIÊU NGHIÊN CỨU
% ==============================================================================

\section{Đặt vấn đề và Bối cảnh nghiên cứu}\label{sec:ch01_dat_van_de}

Phần văn bản chính sử dụng font Times New Roman cỡ chữ 13bp, giãn dòng 1.2 theo đúng Quyết định 1799 của Trường Đại học Cần Thơ \cite{ctu_qd1799}. Dòng đầu tiên của mỗi đoạn văn tự động thụt vào 1.0 cm và các đoạn cách nhau 6bp.

\subsection{Mục tiêu nghiên cứu}

Theo quy định trình bày, đề tài cần làm rõ mục tiêu tổng quát và các mục tiêu cụ thể \cite{sample_article2024}:

\subsubsection{Mục tiêu tổng quát}

Trình bày mục tiêu bao quát mà đề tài hoặc báo cáo khoa học hướng đến giải quyết trong phạm vi nghiên cứu đã xác định.

\subsubsection{Mục tiêu cụ thể}

Các gạch đầu dòng hoặc đánh số liệt kê sử dụng tùy chọn \texttt{[leftmargin=*]} để thẳng hàng đầu dòng 1.0 cm với các tiểu mục và không chèn thêm khoảng trống giữa các dòng liệt kê:
\begin{enumerate}[leftmargin=*]
  \item \textbf{Mục tiêu 1:} Hệ thống hóa cơ sở lý thuyết và các công trình nghiên cứu liên quan đến bài toán.
  \item \textbf{Mục tiêu 2:} Đề xuất quy trình, mô hình hoặc phương pháp giải quyết vấn đề nghiên cứu.
  \item \textbf{Mục tiêu 3:} Triển khai thực nghiệm, đánh giá kết quả định lượng và rút ra các nhận định khoa học.
\end{enumerate}

\section{Phạm vi và Đóng góp của đề tài}\label{sec:ch01_pham_vi}

Các nội dung liệt kê dạng gạch đầu dòng dùng môi trường \texttt{itemize}:
\begin{itemize}[leftmargin=*]
  \item \textbf{Phạm vi dữ liệu và thời gian:} Mô tả giới hạn tập dữ liệu, đối tượng khảo sát hoặc phạm vi thử nghiệm.
  \item \textbf{Đóng góp dự kiến:} Nêu rõ giá trị khoa học và khả năng ứng dụng thực tiễn của kết quả đạt được.
\end{itemize}
